% !TeX root = ../main.tex
% ============================================================
%  appendix/app_data.tex
% ============================================================

\section{Instance Data and Parameter Values}
\label{app:data}

The committed benchmark data are the JSON files under
\path{data/instances/base}. The generator entrypoint is
\path{data/generator/instance.py} (\texttt{generate\_instance}/
\texttt{write\_instance}), invoked via
\texttt{python -m data.generator benchmark}; sampling distributions
and bounds are implemented in \path{data/generator/tautenhain_data.py},
with the instance-tree layout documented in \path{data/README.md}.

\subsection{Benchmark Scope}
\label{app:data:scope}

The current base benchmark contains 12 instances:
\[
  \{\mathrm{I17},\mathrm{I25}\}
  \times
  \{\mathrm{M4},\mathrm{M8}\}
  \times
  \{\mathrm{T3},\mathrm{T6},\mathrm{T12}\}.
\]
The \texttt{STD} prefix is the only prefix generated by the current
benchmark script. The eight T3 and T6 instances are solved to validation; the
four T12 instances are generated but not solved, being intractable within the
compute budget (\Cref{sec:exp:scalability}).

\begin{table}[ht]
\centering
\caption{Implemented benchmark-generation settings.}
\label{tab:app_params}
\renewcommand{\arraystretch}{1.2}
\begin{tabularx}{\linewidth}{@{}lX@{}}
  \toprule
  \textbf{Setting} & \textbf{Implemented value} \\
  \midrule
  Suppliers & 3 \\
  Factories & 3 \\
  Retailers & 4 \\
  Warehouses & 3 for I17; 9 for I25 \\
  Airports & 2 for I17; 4 for I25 \\
  Seaports & 2 \\
  Raw materials and products
    & M4: 2 raw materials and 1 product; M8: 4 raw materials and
      2 products \\
  Production technologies & 3 \\
  Remanufacturing technologies & 3 \\
  Planning periods & T3, T6, or T12 \\
  Coordinate span & \SI{750}{\kilo\metre} square \\
  Tax-index domain & $\taxindex\in\{0,\ldots,20\}$ \\
  Effective tax rate & $\tau=10\,\taxindex$ USD/t\,CO\textsubscript{2}e, i.e.\ $\tau\in[0,200]$ USD/t \\
  \bottomrule
\end{tabularx}
\end{table}

\subsection{Carbon-Tax Range Justification}
\label{app:data:tax_range}

The interval $[0,200]$ USD/t\,CO\textsubscript{2}e was chosen to span
the range from no tax to levels consistent with ambitious
decarbonisation trajectories. The current global average price stands
at approximately \$19/t \citep{worldbank2025}, up from about \$10/t a
decade earlier; the range recommended to meet the Paris Agreement
targets is \$50--\$100/t by 2030 (\$40--\$80/t by 2020)
\citep{stiglitzstern2017}; the highest carbon tax currently implemented
anywhere in the world, Sweden's, stands at approximately \$130--\$145/t
\citep{oecd2024sweden}; and IMF analyses point to shadow prices of
\$150--\$200+/t under scenarios consistent with net-zero emissions by
2050 \citep{imf2024netzero}. The base ceiling $\taxindexub=20$ (200 USD/t)
therefore covers the full policy-relevant spectrum, from today's market
average to the upper end of what net-zero-consistent pathways require.

\subsection{Seed Protocol}
\label{app:data:seeds}

Benchmark seeds are deterministic functions of the benchmark prefix
and suffix. The implementation hashes
\path{seed_base:prefix:suffix} with SHA-256 and uses the first
eight hexadecimal digits as the integer seed. The generated seed is
stored in each instance under \path{metadata.benchmark.seed}.

\subsection{Data Availability}
\label{app:data:availability}

All instance files, generation scripts, solver scripts, tests, and
saved result artifacts are part of the accompanying repository. The
reported numerical results should be regenerated with:
\begin{verbatim}
python -m pytest -q
python tools/solver/solve_batch.py data/instances/base/STD_I17_M4_T3.json
python -m tools.experiments.sensitivity analyze-all
\end{verbatim}
The final command regenerates every derived table, figure, and report under
\path{results/analysis} from the saved solver outputs, without re-solving.
